% META: {"id":"1SPE-PROBCOND-FR-R1","chapitre":"1SPE-PROBA-COND","type_objet":"remediation","status":"approved","prerequis_testes":["R1"],"origin":"generated","status_history":["generated","audited","accepted","approved"],"release_acceptance":"RELEASE_OWNER_BATCH_ACCEPTANCE","acceptance_closure_digest":"sha256:9b3ccf9a81c5520fbb7e03b7b2d3e7bf2057a02834b6d908bf3be5f49f3b553f","acceptance_receipt_digest":"sha256:4fad86f0282e0fa4e0419cca1ad6379ae7af5a280c04a4a90880f90a1c044242"}

%---------------------------------------------------------------
% FICHE DE REMISE A NIVEAU --- R1 : Probabilités de base
%---------------------------------------------------------------

\subsection*{Rappel --- Probabilités de base}

\textbf{Vocabulaire :} un \textbf{univers} $\Omega$ est l'ensemble de toutes les issues possibles. Un \textbf{événement} est une partie de $\Omega$.

\textbf{Formules fondamentales :}
\[
  P(\overline{A}) = 1 - P(A) \qquad P(A \cup B) = P(A) + P(B) - P(A \cap B)
\]

\begin{exercice}{1SPE-PROBCOND-FR-R1-EX1}{1}{5}
On donne $P(A) = \dfrac{3}{8}$ et $P(B) = \dfrac{1}{4}$ avec $A \cap B = \varnothing$.
\begin{enumerate}
  \item Calculer $P(\overline{A})$.
  \item Calculer $P(A \cup B)$.
\end{enumerate}
\end{exercice}

% BEGIN-VERIFY
% from sympy import Rational
% P_A = Rational(3, 8)
% P_B = Rational(1, 4)
% assert 1 - P_A == Rational(5, 8)
% assert P_A + P_B == Rational(5, 8)
% END-VERIFY

\begin{corrige}{1SPE-PROBCOND-FR-R1-EX1}
\textbf{1.} $P(\overline{A}) = 1 - P(A) = 1 - \frac{3}{8} = \frac{5}{8}$.

\textbf{2.} Comme $A \cap B = \varnothing$, $P(A \cup B) = P(A) + P(B) = \frac{3}{8} + \frac{1}{4} = \frac{3}{8} + \frac{2}{8} = \frac{5}{8}$.
\end{corrige}
